% =====================================================================
% VERIFICA DI MATEMATICA - 2BS
% File A e B: 60 minuti e 30 punti per fila.
% Punteggi per fila: 4 + 8 + 10 + 3 + 5 = 30.
% Compilare almeno due volte con pdfLaTeX.
% Per mostrare le soluzioni: decommentare \printanswers.
% =====================================================================

\documentclass[11pt,a4paper,addpoints]{exam}

% ---------- Lingua e codifica ----------
\usepackage[utf8]{inputenc}
\usepackage[T1]{fontenc}
\usepackage[italian]{babel}
\usepackage{lmodern}
\usepackage{microtype}

% ---------- Matematica e impaginazione ----------
\usepackage{amsmath,amssymb}
\usepackage{xcolor}
\usepackage{array,tabularx}
\usepackage[
  a4paper,
  top=2cm,
  bottom=2.2cm,
  left=2cm,
  right=2.2cm,
  headheight=30pt,
  headsep=12pt
]{geometry}

% ---------- Soluzioni ----------
%\printanswers
\ifdefined\CORRETTORE
  \printanswers
\fi
\cancelspace
\unframedsolutions
\renewcommand{\solutiontitle}{%
  \noindent\textbf{Soluzione.}\enspace
}
\SolutionEmphasis{\color{blue!60!black}}

% ---------- Dati della prova ----------
\newcommand{\materia}{Matematica}
\newcommand{\classe}{2BS}
\newcommand{\indirizzo}{Liceo Scientifico}
\newcommand{\tipoprova}{Verifica di matematica}
\newcommand{\argomento}{Frazioni algebriche ed equazioni frazionarie}
\newcommand{\durata}{60 minuti}
\newcommand{\annoscolastico}{2026/2027}
\newcommand{\docente}{Prof. Rubens Longhi}
\newcommand{\fila}{A}

\newcommand{\indicazioni}{%
  Svolgi sul foglio a protocollo e mostra i passaggi.
  Indica sempre le condizioni di esistenza prima di semplificare.
  Per le equazioni scrivi l'insieme delle soluzioni.
  Lavora in $\mathbb{R}$, senza calcolatrice.
}

% ---------- Macro matematiche ----------
\newcommand{\R}{\mathbb{R}}

% ---------- Intestazioni e piè di pagina ----------
\pagestyle{headandfoot}
\firstpageheader
  {%
    \scriptsize\textsc{IIS David Maria Turoldo}\\[-1pt]
    \scriptsize Via Ronco 11 -- 24019 Zogno (BG)
  }
  {}
  {%
    \scriptsize A.S.\ \annoscolastico\\[-1pt]
    \scriptsize \docente
  }
\firstpageheadrule
\runningheader
  {\small\materia\ -- \classe}
  {\small\tipoprova}
  {\small Fila \fila}
\runningheadrule
\newcommand{\paginefila}{\pageref{fine@fila@\fila}}
\firstpagefooter{}{\small Pagina \thepage\ di \paginefila}{}
\runningfooter{}{\small Pagina \thepage\ di \paginefila}{}

% ---------- Obiettivi e formato degli esercizi ----------
\newcommand{\obiettivocorrente}{}
\newcommand{\obiettivo}[1]{%
  \gdef\obiettivocorrente{#1}%
}
\newcommand{\stampaobiettivo}{%
  \ifx\obiettivocorrente\empty
  \else
    {\footnotesize\color{black!55}%
      \textit{Obiettivo: \obiettivocorrente}}%
  \fi
}
\qformat{%
  \textbf{\thequestion.}\enspace
  \stampaobiettivo
  \hfill
  {\small(\totalpoints\ \points)}%
  \gdef\obiettivocorrente{}%
}
\renewcommand{\questionlabel}{}
\renewcommand{\partlabel}{\thequestion\thepartno)}
\renewcommand{\subpartlabel}{\thesubpart.}
\pointpoints{punto}{punti}
\renewcommand{\partshook}{%
  \setlength{\itemsep}{4pt}%
  \setlength{\rightmargin}{0.5cm}%
}

% ---------- Numerazione indipendente per ciascuna fila ----------
\newcounter{offsetfila}
\newcounter{totalqfila}
\renewcommand{\thequestion}{%
  \number\numexpr
    \value{question}-\value{offsetfila}%
  \relax
}
\newenvironment{esercizi}
  {%
    \begin{questions}
    \setcounter{question}{\value{offsetfila}}%
  }
  {\end{questions}}

% ---------- Esercizi non spezzati nella versione studenti ----------
\newenvironment{esercizio}
  {%
    \par
    \ifprintanswers\else\begin{samepage}\fi
  }
  {%
    \ifprintanswers\else\end{samepage}\fi
    \par\vspace{3pt}
  }

% ---------- Tabella dei punteggi ----------
\hqword{\textbf{Esercizio}}
\hpword{\textbf{Punti massimi}}
\hsword{\textbf{Punti ottenuti}}
\htword{\textbf{Totale}}

% ---------- Inizio e fine di ciascuna fila ----------
\newcommand{\iniziofila}[1]{%
  \clearpage
  \renewcommand{\fila}{#1}%
  \setcounter{page}{1}%
  \setcounter{offsetfila}{\value{totalqfila}}%
  \begin{center}
    {\small\bfseries
      \tipoprova
      \ifprintanswers
        -- Soluzioni e criteri di correzione
      \fi
    }\\[3pt]
    {\Large \argomento}\\[2pt]
    {\small
      \indirizzo\ -- Classe \classe\
      -- Durata: \durata\
      -- \textbf{Fila #1}
    }
  \end{center}
  \vspace{4pt}
  \noindent
  \begin{tabularx}{\textwidth}{@{}l X l p{3cm}@{}}
    \textbf{Nome e Cognome:} & \hrulefill
    & \textbf{Data:} & \hrulefill\\[8pt]
  \end{tabularx}
  \noindent
  \fbox{%
    \parbox{\dimexpr\textwidth-2\fboxsep-2\fboxrule}{%
      \small
      \textbf{Indicazioni.} \indicazioni\space
      Punteggio totale: \pointsinrange{fila#1} punti.
    }%
  }
  \medskip
  \begingradingrange{fila#1}%
}
\newcommand{\finefila}{%
  \setcounter{totalqfila}{\value{question}}%
  \endgradingrange{fila\fila}%
  \ifprintanswers
  \else
    \vfill
    \begin{center}
      \small
      \partialgradetable{fila\fila}[h][questions]
    \end{center}
    \vspace{6pt}
    \noindent
    {\small
      \textbf{Voto:}\ \rule{3cm}{0.4pt}
      \hfill
      \textbf{Firma del docente:}\ \rule{5cm}{0.4pt}
    }%
  \fi
  \label{fine@fila@\fila}%
}

% =====================================================================
% DOCUMENTO
% =====================================================================
\begin{document}
\linespread{1.05}\selectfont
\ifprintanswers\small\fi

% =====================================================================
% FILA A - 30 punti
% =====================================================================
\iniziofila{A}
\begin{esercizi}

% ---------------- ESERCIZIO 1: 4 punti ----------------
\begin{esercizio}
\obiettivo{ragionare su equivalenza, C.E. e soluzioni}
\question
Stabilisci se ciascuna affermazione è vera o falsa.
Giustifica ogni risposta; per le affermazioni false
puoi anche fornire un controesempio.

\begin{parts}
\part[1]
La frazione $\dfrac{x^2-4}{2-x}$ è equivalente
all'espressione $x+2$ per tutti i valori di $x$
per cui entrambe sono definite.
\begin{solution}
\textbf{Falsa.} Per $x\ne2$:
\[
\frac{x^2-4}{2-x}
=\frac{(x-2)(x+2)}{-(x-2)}=-(x+2).
\]
Per esempio, per $x=0$ si ottengono $-2$ e $2$.
\end{solution}

\part[1]
L'espressione
$\dfrac{x^2-4}{x-2}:\dfrac{x+2}{x-1}$
ha come C.E. soltanto $x\ne1$ e $x\ne2$.
\begin{solution}
\textbf{Falsa.} I denominatori richiedono $x\ne2,1$;
il divisore deve essere non nullo, quindi occorre anche $x\ne-2$.
Le C.E. complete sono $\boxed{x\ne-2,1,2}$,
anche se l'espressione si semplifica in $x-1$.
\end{solution}

\part[1]
Per ogni $a\in\R$ con $a\ne0$, l'equazione
$\dfrac{x-1}{a}=\dfrac{a}{x-1}$
ha esattamente due soluzioni reali distinte.
\begin{solution}
\textbf{Vera.} Le C.E. sono $x\ne1$ e il parametro è non nullo.
\[
(x-1)^2=a^2
\iff (x-1-a)(x-1+a)=0,
\qquad S=\{1-a,1+a\}.
\]
Poiché $a\ne0$, le due soluzioni sono distinte ed entrambe ammesse.
\end{solution}

\part[1]
L'equazione
$\dfrac{x+2}{x-2}-\dfrac{8x}{x^2-4}=\dfrac{x-2}{x+2}$
è indeterminata e ha insieme delle soluzioni $\R\setminus\{-2,2\}$.
\begin{solution}
\textbf{Vera.} Le C.E. sono $x\ne-2,2$.
Moltiplicando per $(x-2)(x+2)$ si ottiene l'identità
\[
(x+2)^2-8x=(x-2)^2.
\]
Le soluzioni sono tutti e soli i valori ammessi dalle C.E.
\end{solution}
\end{parts}
\begin{solution}
\textit{Criteri per ciascuna risposta:}
scelta 0,5 punti; giustificazione corretta 0,5 punti.
\end{solution}
\end{esercizio}

% ---------------- ESERCIZIO 2: 8 punti ----------------
\begin{esercizio}
\obiettivo{operare con frazioni algebriche e potenze negative}
\question
Determina le C.E. e semplifica le seguenti espressioni.
\begin{parts}
\part[4]
\[
\left[
\frac{x}{x^2-5x+6}+\frac1{x-2}-\frac1{x-3}
\right]:\frac{2x-2}{x^2-4}.
\]
\begin{solution}
Poiché $x^2-5x+6=(x-2)(x-3)$ e $x^2-4=(x-2)(x+2)$,
i denominatori e il divisore impongono
$\boxed{\text{C.E.: }x\ne-2,1,2,3}$.
\begin{align*}
\left[
\frac{x}{(x-2)(x-3)}+\frac1{x-2}-\frac1{x-3}
\right]
&=\frac{x+(x-3)-(x-2)}{(x-2)(x-3)}
=\frac{x-1}{(x-2)(x-3)},\\
\frac{x-1}{(x-2)(x-3)}
\cdot\frac{(x-2)(x+2)}{2(x-1)}
&=\boxed{\frac{x+2}{2(x-3)}}.
\end{align*}
Restano valide tutte le C.E. iniziali.
\textit{Criteri:} C.E. 1 punto; somma 1 punto;
divisione e semplificazioni 1,5 punti; risultato 0,5 punti.
\end{solution}

\part[4]
\[
\left[
\left(\frac{a^{-1}-3a^{-2}}{a^{-1}+3a^{-2}}\right)^2
+2\frac{a^{-1}-3a^{-2}}{a^{-1}+3a^{-2}}+1
\right]:\frac{2a^2}{a^3+6a^2+9a}.
\]
\begin{solution}
Le potenze negative richiedono $a\ne0$.
Inoltre $a^{-1}+3a^{-2}=(a+3)/a^2$ e
$a^3+6a^2+9a=a(a+3)^2$.
Il divisore deve essere non nullo:
$\boxed{\text{C.E.: }a\ne0,-3}$.
\[
\frac{a^{-1}-3a^{-2}}{a^{-1}+3a^{-2}}
=\frac{a-3}{a+3}.
\]
La parentesi è il quadrato di un binomio:
\[
\left(\frac{a-3}{a+3}\right)^2
+2\frac{a-3}{a+3}+1
=\left(\frac{a-3}{a+3}+1\right)^2
=\frac{4a^2}{(a+3)^2}.
\]
Pertanto
\[
\frac{4a^2}{(a+3)^2}\cdot\frac{a(a+3)^2}{2a^2}
=\boxed{2a},
\]
mantenendo le C.E. iniziali.
\textit{Criteri:} C.E. 1 punto; potenze negative 1 punto;
parentesi quadra 1 punto; divisione e risultato 1 punto.
\end{solution}
\end{parts}
\end{esercizio}

\ifprintanswers\else\newpage\fi

% ---------------- ESERCIZIO 3: 10 punti ----------------
\begin{esercizio}
\obiettivo{risolvere equazioni e controllare l'ammissibilità}
\question
Risolvi le seguenti equazioni, indicando le C.E.
e l'insieme delle soluzioni.
\begin{parts}
\part[3]
\[
\frac{x^2+2x-9}{x^2-9}=\frac1{x-3}.
\]
\begin{solution}
$\boxed{\text{C.E.: }x\ne-3,3}$.
Moltiplicando per $(x-3)(x+3)$ nel dominio ammesso:
\[
x^2+2x-9=x+3
\iff x^2+x-12=0
\iff (x+4)(x-3)=0.
\]
I candidati sono $x=-4$ e $x=3$; il secondo è escluso.
Dunque $\boxed{S=\{-4\}}$.
\textit{Criteri:} C.E. 1 punto;
riduzione e scomposizione 1 punto;
selezione delle soluzioni 1 punto.
\end{solution}

\part[3]
\[
-\frac2x+\frac5{x+1}=\frac{x}{x+1}.
\]
\begin{solution}
$\boxed{\text{C.E.: }x\ne0,-1}$.
Moltiplicando per $x(x+1)$:
\[
-2(x+1)+5x=x^2
\iff x^2-3x+2=0
\iff (x-1)(x-2)=0.
\]
Entrambi i valori sono ammessi: $\boxed{S=\{1,2\}}$.
\textit{Criteri:} C.E. 0,5 punti; riduzione 1 punto;
scomposizione 1 punto; insieme delle soluzioni 0,5 punti.
\end{solution}

\part[4]
\[
\left(\frac{2x+5}{x^2-9}+\frac1{3-x}\right)
:\left(\frac2{9-x^2}\right)
=\frac{15-5x}{2x-6}.
\]
\begin{solution}
I denominatori richiedono $x\ne-3,3$.
Il divisore non è mai nullo dove esiste:
\[
\boxed{\text{C.E.: }x\ne-3,3}.
\]
\[
\frac{2x+5}{(x-3)(x+3)}-\frac1{x-3}
=\frac{x+2}{(x-3)(x+3)}.
\]
Poiché $9-x^2=-(x-3)(x+3)$ e
$15-5x=-5(x-3)$, nel dominio ammesso:
\[
-\frac{x+2}{2}=-\frac52
\iff x=3.
\]
Il candidato è escluso dalle C.E.:
$\boxed{S=\varnothing}$.
\textit{Criteri:} C.E. 1 punto; somma e segni 1 punto;
divisione e riduzione 1 punto; esclusione e conclusione 1 punto.
\end{solution}
\end{parts}
\end{esercizio}

% ---------------- ESERCIZIO 4: 3 punti ----------------
\begin{esercizio}
\obiettivo{ricavare una formula inversa}
\question[3]
Un veicolo percorre due tratti di uguale lunghezza con velocità
costanti positive $v_1$ e $v_2$.
La velocità media sull'intero percorso è
\[
v_m=\frac{2v_1v_2}{v_1+v_2}.
\]
Ricava $v_2$ in funzione di $v_m$ e $v_1$.
Sapendo che $v_m>0$ e $v_1>0$, indica quale condizione
devono soddisfare affinché il valore ricavato di $v_2$
esista e sia positivo.
\begin{solution}
Essendo $v_1+v_2>0$:
\[
v_m(v_1+v_2)=2v_1v_2
\iff v_mv_1=(2v_1-v_m)v_2.
\]
La formula inversa è
\[
\boxed{v_2=\frac{v_mv_1}{2v_1-v_m}},
\qquad 2v_1-v_m\ne0.
\]
Il numeratore è positivo. Affinché la formula esista e
fornisca un valore positivo occorre e basta che
\[
\boxed{0<v_m<2v_1}.
\]
Se $v_m=2v_1$, il denominatore è zero e la formula non è definita;
nell'equazione si avrebbe $v_mv_1=0$, impossibile con i dati positivi.
\textit{Criteri:} passaggi e formula inversa 2 punti;
condizione di esistenza e positività 1 punto.
\end{solution}
\end{esercizio}

% ---------------- ESERCIZIO 5: 5 punti ----------------
\begin{esercizio}
\obiettivo{interpretare e confrontare modelli frazionari}
\question
Le temperature di due componenti di una macchina, in gradi Celsius,
sono descritte dai modelli
\[
T_1(t)=\frac{60(t+1)}{t+3},
\qquad T_2(t)=\frac{30(2t+1)}{t+2},
\]
dove $t\ge0$ è il numero di minuti trascorsi dall'accensione.
\begin{parts}
\part[1]
Quali sono le temperature dei due componenti
all'istante dell'accensione?
\begin{solution}
$T_1(0)=60/3=\boxed{20\,^\circ\mathrm{C}}$;
$T_2(0)=30/2=\boxed{15\,^\circ\mathrm{C}}$.
\textit{Criteri:} 0,5 punti per ciascuna temperatura.
\end{solution}

\part[2]
Dopo quanti minuti il primo componente raggiunge
la temperatura di $40\,^\circ\mathrm{C}$?
Imposta e risolvi un'equazione.
\begin{solution}
Per $t\ge0$, il denominatore è non nullo.
\[
\frac{60(t+1)}{t+3}=40
\iff60t+60=40t+120
\iff t=3.
\]
Il valore è ammissibile: $\boxed{3\text{ minuti}}$.
\textit{Criteri:} equazione e dominio 1 punto;
risoluzione e risposta 1 punto.
\end{solution}

\part[2]
Dopo quanti minuti i due componenti hanno la stessa temperatura?
Qual è tale temperatura? Imposta e risolvi un'equazione.
\begin{solution}
Per $t\ge0$, entrambi i denominatori sono non nulli.
\begin{align*}
\frac{60(t+1)}{t+3}&=\frac{30(2t+1)}{t+2}\\
2(t+1)(t+2)&=(2t+1)(t+3)\\
2t^2+6t+4&=2t^2+7t+3\\
t&=1.
\end{align*}
Il valore è ammissibile:
$T_1(1)=T_2(1)=30\,^\circ\mathrm{C}$.
Quindi $\boxed{1\text{ minuto}}$ e $\boxed{30\,^\circ\mathrm{C}}$.
\textit{Criteri:} equazione e passaggi 1 punto;
tempo ammissibile 0,5 punti; temperatura 0,5 punti.
\end{solution}
\end{parts}
\end{esercizio}
\end{esercizi}
\finefila

% =====================================================================
% FILA B - 30 punti
% =====================================================================
\iniziofila{B}
\begin{esercizi}

% ---------------- ESERCIZIO 1: 4 punti ----------------
\begin{esercizio}
\obiettivo{ragionare su equivalenza, C.E. e soluzioni}
\question
Stabilisci se ciascuna affermazione è vera o falsa.
Giustifica ogni risposta; per le affermazioni false
puoi anche fornire un controesempio.
\begin{parts}
\part[1]
La frazione $\dfrac{x^2-9}{3-x}$ è equivalente
all'espressione $x+3$ per tutti i valori di $x$
per cui entrambe sono definite.
\begin{solution}
\textbf{Falsa.} Per $x\ne3$:
\[
\frac{x^2-9}{3-x}
=\frac{(x-3)(x+3)}{-(x-3)}=-(x+3).
\]
Per esempio, per $x=0$ si ottengono $-3$ e $3$.
\end{solution}

\part[1]
L'espressione
$\dfrac{x^2-9}{x-3}:\dfrac{x+3}{x-2}$
ha come C.E. soltanto $x\ne2$ e $x\ne3$.
\begin{solution}
\textbf{Falsa.} I denominatori richiedono $x\ne3,2$;
il divisore deve essere non nullo, quindi occorre anche $x\ne-3$.
Le C.E. complete sono $\boxed{x\ne-3,2,3}$,
anche se l'espressione si semplifica in $x-2$.
\end{solution}

\part[1]
Per ogni $a\in\R$ con $a\ne0$, l'equazione
$\dfrac{x-2}{a}=\dfrac{a}{x-2}$
ha esattamente due soluzioni reali distinte.
\begin{solution}
\textbf{Vera.} Le C.E. sono $x\ne2$ e il parametro è non nullo.
\[
(x-2)^2=a^2
\iff(x-2-a)(x-2+a)=0,
\qquad S=\{2-a,2+a\}.
\]
Poiché $a\ne0$, le due soluzioni sono distinte ed entrambe ammesse.
\end{solution}

\part[1]
L'equazione
$\dfrac{x+3}{x-3}-\dfrac{12x}{x^2-9}=\dfrac{x-3}{x+3}$
è indeterminata e ha insieme delle soluzioni $\R\setminus\{-3,3\}$.
\begin{solution}
\textbf{Vera.} Le C.E. sono $x\ne-3,3$.
Moltiplicando per $(x-3)(x+3)$ si ottiene l'identità
\[
(x+3)^2-12x=(x-3)^2.
\]
Le soluzioni sono tutti e soli i valori ammessi dalle C.E.
\end{solution}
\end{parts}
\begin{solution}
\textit{Criteri per ciascuna risposta:}
scelta 0,5 punti; giustificazione corretta 0,5 punti.
\end{solution}
\end{esercizio}

% ---------------- ESERCIZIO 2: 8 punti ----------------
\begin{esercizio}
\obiettivo{operare con frazioni algebriche e potenze negative}
\question
Determina le C.E. e semplifica le seguenti espressioni.
\begin{parts}
\part[4]
\[
\left[
\frac{x}{x^2-8x+15}+\frac1{x-3}-\frac1{x-5}
\right]:\frac{3x-6}{x^2-9}.
\]
\begin{solution}
Poiché $x^2-8x+15=(x-3)(x-5)$ e $x^2-9=(x-3)(x+3)$,
i denominatori e il divisore impongono
$\boxed{\text{C.E.: }x\ne-3,2,3,5}$.
\begin{align*}
\left[
\frac{x}{(x-3)(x-5)}+\frac1{x-3}-\frac1{x-5}
\right]
&=\frac{x+(x-5)-(x-3)}{(x-3)(x-5)}
=\frac{x-2}{(x-3)(x-5)},\\
\frac{x-2}{(x-3)(x-5)}
\cdot\frac{(x-3)(x+3)}{3(x-2)}
&=\boxed{\frac{x+3}{3(x-5)}}.
\end{align*}
Restano valide tutte le C.E. iniziali.
\textit{Criteri:} C.E. 1 punto; somma 1 punto;
divisione e semplificazioni 1,5 punti; risultato 0,5 punti.
\end{solution}

\part[4]
\[
\left[
\left(\frac{a^{-1}-4a^{-2}}{a^{-1}+4a^{-2}}\right)^2
+2\frac{a^{-1}-4a^{-2}}{a^{-1}+4a^{-2}}+1
\right]:\frac{3a^2}{a^3+8a^2+16a}.
\]
\begin{solution}
Le potenze negative richiedono $a\ne0$.
Inoltre $a^{-1}+4a^{-2}=(a+4)/a^2$ e
$a^3+8a^2+16a=a(a+4)^2$.
Il divisore deve essere non nullo:
$\boxed{\text{C.E.: }a\ne0,-4}$.
\[
\frac{a^{-1}-4a^{-2}}{a^{-1}+4a^{-2}}
=\frac{a-4}{a+4}.
\]
La parentesi è il quadrato di un binomio:
\[
\left(\frac{a-4}{a+4}\right)^2
+2\frac{a-4}{a+4}+1
=\left(\frac{a-4}{a+4}+1\right)^2
=\frac{4a^2}{(a+4)^2}.
\]
Pertanto
\[
\frac{4a^2}{(a+4)^2}\cdot\frac{a(a+4)^2}{3a^2}
=\boxed{\frac{4a}{3}},
\]
mantenendo le C.E. iniziali.
\textit{Criteri:} C.E. 1 punto; potenze negative 1 punto;
parentesi quadra 1 punto; divisione e risultato 1 punto.
\end{solution}
\end{parts}
\end{esercizio}

\ifprintanswers\else\newpage\fi

% ---------------- ESERCIZIO 3: 10 punti ----------------
\begin{esercizio}
\obiettivo{risolvere equazioni e controllare l'ammissibilità}
\question
Risolvi le seguenti equazioni, indicando le C.E.
e l'insieme delle soluzioni.
\begin{parts}
\part[3]
\[
\frac{x^2+2x-16}{x^2-16}=\frac1{x-4}.
\]
\begin{solution}
$\boxed{\text{C.E.: }x\ne-4,4}$.
Moltiplicando per $(x-4)(x+4)$ nel dominio ammesso:
\[
x^2+2x-16=x+4
\iff x^2+x-20=0
\iff(x+5)(x-4)=0.
\]
I candidati sono $x=-5$ e $x=4$; il secondo è escluso.
Dunque $\boxed{S=\{-5\}}$.
\textit{Criteri:} C.E. 1 punto;
riduzione e scomposizione 1 punto;
selezione delle soluzioni 1 punto.
\end{solution}

\part[3]
\[
-\frac3x+\frac8{x+2}=\frac{x}{x+2}.
\]
\begin{solution}
$\boxed{\text{C.E.: }x\ne0,-2}$.
Moltiplicando per $x(x+2)$:
\[
-3(x+2)+8x=x^2
\iff x^2-5x+6=0
\iff(x-2)(x-3)=0.
\]
Entrambi i valori sono ammessi: $\boxed{S=\{2,3\}}$.
\textit{Criteri:} C.E. 0,5 punti; riduzione 1 punto;
scomposizione 1 punto; insieme delle soluzioni 0,5 punti.
\end{solution}

\part[4]
\[
\left(\frac{2x+7}{x^2-16}+\frac1{4-x}\right)
:\left(\frac2{16-x^2}\right)
=\frac{28-7x}{2x-8}.
\]
\begin{solution}
I denominatori richiedono $x\ne-4,4$.
Il divisore non è mai nullo dove esiste:
\[
\boxed{\text{C.E.: }x\ne-4,4}.
\]
\[
\frac{2x+7}{(x-4)(x+4)}-\frac1{x-4}
=\frac{x+3}{(x-4)(x+4)}.
\]
Poiché $16-x^2=-(x-4)(x+4)$ e
$28-7x=-7(x-4)$, nel dominio ammesso:
\[
-\frac{x+3}{2}=-\frac72
\iff x=4.
\]
Il candidato è escluso dalle C.E.:
$\boxed{S=\varnothing}$.
\textit{Criteri:} C.E. 1 punto; somma e segni 1 punto;
divisione e riduzione 1 punto; esclusione e conclusione 1 punto.
\end{solution}
\end{parts}
\end{esercizio}

% ---------------- ESERCIZIO 4: 3 punti ----------------
\begin{esercizio}
\obiettivo{ricavare una formula inversa}
\question[3]
Un veicolo percorre due tratti di uguale lunghezza con velocità
costanti positive $v_1$ e $v_2$.
La velocità media sull'intero percorso è
\[
v_m=\frac{2v_1v_2}{v_1+v_2}.
\]
Ricava $v_1$ in funzione di $v_m$ e $v_2$.
Sapendo che $v_m>0$ e $v_2>0$, indica quale condizione
devono soddisfare affinché il valore ricavato di $v_1$
esista e sia positivo.
\begin{solution}
Essendo $v_1+v_2>0$:
\[
v_m(v_1+v_2)=2v_1v_2
\iff v_mv_2=(2v_2-v_m)v_1.
\]
La formula inversa è
\[
\boxed{v_1=\frac{v_mv_2}{2v_2-v_m}},
\qquad 2v_2-v_m\ne0.
\]
Il numeratore è positivo. Affinché la formula esista e
fornisca un valore positivo occorre e basta che
\[
\boxed{0<v_m<2v_2}.
\]
Se $v_m=2v_2$, il denominatore è zero e la formula non è definita;
nell'equazione si avrebbe $v_mv_2=0$, impossibile con i dati positivi.
\textit{Criteri:} passaggi e formula inversa 2 punti;
condizione di esistenza e positività 1 punto.
\end{solution}
\end{esercizio}

% ---------------- ESERCIZIO 5: 5 punti ----------------
\begin{esercizio}
\obiettivo{interpretare e confrontare modelli frazionari}
\question
Le temperature di due componenti di una macchina, in gradi Celsius,
sono descritte dai modelli
\[
T_1(t)=\frac{120(t+1)}{t+4},
\qquad T_2(t)=\frac{60(2t+1)}{t+3},
\]
dove $t\ge0$ è il numero di minuti trascorsi dall'accensione.
\begin{parts}
\part[1]
Quali sono le temperature dei due componenti
all'istante dell'accensione?
\begin{solution}
$T_1(0)=120/4=\boxed{30\,^\circ\mathrm{C}}$;
$T_2(0)=60/3=\boxed{20\,^\circ\mathrm{C}}$.
\textit{Criteri:} 0,5 punti per ciascuna temperatura.
\end{solution}

\part[2]
Dopo quanti minuti il primo componente raggiunge
la temperatura di $80\,^\circ\mathrm{C}$?
Imposta e risolvi un'equazione.
\begin{solution}
Per $t\ge0$, il denominatore è non nullo.
\[
\frac{120(t+1)}{t+4}=80
\iff120t+120=80t+320
\iff t=5.
\]
Il valore è ammissibile: $\boxed{5\text{ minuti}}$.
\textit{Criteri:} equazione e dominio 1 punto;
risoluzione e risposta 1 punto.
\end{solution}

\part[2]
Dopo quanti minuti i due componenti hanno la stessa temperatura?
Qual è tale temperatura? Imposta e risolvi un'equazione.
\begin{solution}
Per $t\ge0$, entrambi i denominatori sono non nulli.
\begin{align*}
\frac{120(t+1)}{t+4}&=\frac{60(2t+1)}{t+3}\\
2(t+1)(t+3)&=(2t+1)(t+4)\\
2t^2+8t+6&=2t^2+9t+4\\
t&=2.
\end{align*}
Il valore è ammissibile:
$T_1(2)=T_2(2)=60\,^\circ\mathrm{C}$.
Quindi $\boxed{2\text{ minuti}}$ e $\boxed{60\,^\circ\mathrm{C}}$.
\textit{Criteri:} equazione e passaggi 1 punto;
tempo ammissibile 0,5 punti; temperatura 0,5 punti.
\end{solution}
\end{parts}
\end{esercizio}
\end{esercizi}
\finefila
\end{document}

